\cleardoublepage

\chapter*{Resumen\markboth{Resumen}{Resumen}}

La robótica es un campo multidisciplinar de la ingeniería centrada en el diseño, construcción y programación de sistemas robotizados con el objetivo de automatizar tareas y facilitar la vida a las personas. Sus aplicaciones son variadas, aunque en muchos casos implican costes elevados que limitan el acceso para el público general. Por ello, gracias al desarrollo de la robótica de bajo coste, es posible acercar estas tecnologías a un mayor número de usuarios.

\vspace{0.2cm}

En el ámbito asistencial, el progresivo envejecimiento de la población y la creciente demanda de cuidados plantea desafíos en cuanto a la disponibilidad de recursos y personal especializado. En este contexto, el presente proyecto aborda el diseño y desarrollo de un robot asistencial de bajo coste orientado al acompañamiento y estimulación de las personas mayores.

\vspace{0.2cm}

La estructura física del robot se ha diseñado mediante herramientas de modelado 3D de acceso libre y se ha fabricado utilizando una impresora 3D estándar. Asimismo, los componentes \textit{hardware} empleados son accesibles y de bajo coste.

\vspace{0.2cm}

El \textit{software} desarrollado integra diferentes tecnologías, como el reconocimiento facial para la autenticación de usuario, la síntesis de voz, y los modelos de lenguaje para la generación de respuestas. Además, se han implementado diversas funcionalidades para el acompañamiento diario, como calendario, recordatorios, notas y juegos de estimulación cognitiva, todas ellas accesibles a través de una interfaz gráfica.

\vspace{0.2cm}

Como unidad de procesamiento del sistema se ha utilizado una placa Raspberry Pi 4, complementada por un ordenador externo encargado de ejecutar los procesos de mayor carga computacional. Gracias a esta arquitectura, se ha conseguido desarrollar un sistema funcional capaz de interactuar con el usuario.

\vspace{0.2cm}

Finalmente, se realizaron diferentes experimentos para la puesta a punto tanto de la parte \textit{software} como la \textit{hardware}, así como evaluar el correcto funcionamiento de las distintas partes de la infraestructura del sistema, y el rendimiento de este en su conjunto.


